\chapter{A Concrete Naming Certificate for $\HeineBorelFiniteNetUp$}
\label{ch:concrete-instances-heine-borel-finite-net-namecert}
\origin{human}

Following Bishop and Bridges' constructive treatment of compact intervals by finite approximation data, the $\HeineBorelFiniteNetUp$ packet records the interval-cover certificate obtained from a located closed interval, a displayed finite net, a positive Lebesgue-number ledger, and the endpoint real-name boundary. It consumes the finite-net interval surface of \autoref{ch:concrete-instances-located-closed-interval-finite-net-namecert}, the compact-cover radius route of \autoref{ch:concrete-instances-lebesgue-number-lemma-namecert}, the compact uniform-modulus route of \autoref{ch:concrete-instances-heine-cantor-namecert}, the terminal real-name seal of \autoref{ch:concrete-instances-real-namecert}, and the dyadic tolerance rows of \autoref{ch:concrete-instances-dyadicratcore-namecert}. The packet is a finite-net interval-cover certificate; it is not a general open-cover compactness theorem, an unrestricted Heine-Borel principle, a selected subcover functional, or an ambient completeness assertion.

\begin{definition}[Heine-Borel finite-net carrier]
\label{def:heine-borel-finite-net-carrier}
A $\HeineBorelFiniteNetUp$ carrier is a finite $\BHist$ packet
$$
H=(I,F,L,C,R,D,E,T,Q,P,N)
$$
with the following displayed rows:
\begin{enumerate}
\item $I$ is the located closed interval source row, exposing the lower endpoint, upper endpoint, and endpoint order data as $\RealUp$ histories.
\item $F$ is the $\LocatedClosedIntervalFiniteNetUp$ finite-net row, listing the displayed mesh scale and finite interval net consumed by the cover request.
\item $L$ is the $\LebesgueNumberLemmaUp$ positive-radius row attached to the displayed finite cover or pointwise radius family.
\item $C$ is the $\HeineCantorUp$ compact-uniformity row that may be read only as a sibling route through the same finite-net and radius data.
\item $R$ is the cover-read row, assigning each listed net point to its displayed cover witness and recording the finite neighborhood test used for nearby interval points.
\item $D$ is the $\DyadicRatCoreUp$ tolerance ledger, recording mesh width, endpoint containment, radius comparison, and finite cover slack checks.
\item $E$ is the terminal $\RealUp$ seal row for the endpoint and cover-boundary reads.
\item $T$ records componentwise $\hsame$ transport for the interval source, finite-net row, Lebesgue-number row, Heine-Cantor row, cover-read row, dyadic ledger, Real seal, replay row, provenance row, and local naming row.
\item $Q$ records displayed $\Cont$ replay from an interval-cover request through $I,F,L,R,D,E$ and, when needed by a uniform-continuity consumer, through the sibling row $C$.
\item $P$ is the $\Pkg$ provenance over $\HeineBorelFiniteNetUp$, $\LocatedClosedIntervalFiniteNetUp$, $\LebesgueNumberLemmaUp$, $\HeineCantorUp$, $\RealUp$, $\DyadicRatCoreUp$, $\BHist$, $\hsame$, $\Cont$, and $\NameCert$.
\item $N$ is the local $\NameCert_{\HeineBorelFiniteNetUp}$ row.
\end{enumerate}
The classifier compares accepted packets only through $I,F,L,C,R,D,E,T,Q,P,N$. It has no coordinate for arbitrary open covers, unrestricted Heine-Borel compactness, selected subcovers, quotient interval representatives, countable choice, host real equality, or ambient $\RealUp$ completeness.
\end{definition}

\begin{theorem}[Heine-Borel finite-net NameCert obligations]
\label{thm:heine-borel-finite-net-namecert-obligations}
For every accepted $\HeineBorelFiniteNetUp$ carrier $H=(I,F,L,C,R,D,E,T,Q,P,N)$, the local $\NameCert_{\HeineBorelFiniteNetUp}$ surface consists exactly of five finite obligations:
\begin{enumerate}
\item carrier admission for the located interval source, finite-net row, Lebesgue-number row, Heine-Cantor sibling row, cover-read row, dyadic ledger, Real seal, transport, replay, provenance, and local naming rows;
\item interval finite-net exposure, requiring the cover request to display $I$ and $F$ before any cover witness is read;
\item radius compatibility, requiring the positive-radius row $L$ and dyadic ledger $D$ to control mesh width, endpoint containment, cover slack, and finite neighborhood tests;
\item sibling uniformity compatibility, requiring the $\HeineCantorUp$ row $C$ to be read only as a sibling route over the same finite-net and radius rows;
\item non-escape, requiring $T,Q,P,N$ to transport, replay, source, and name only the displayed finite packet and to export no arbitrary open-cover compactness, unrestricted Heine-Borel principle, selected subcover, countable-choice row, host real equality, or ambient $\RealUp$ completeness coordinate.
\end{enumerate}
\end{theorem}
\begin{proof}
Carrier admission is projection from \autoref{def:heine-borel-finite-net-carrier}. The tuple displays exactly the interval source, finite net, positive-radius row, sibling uniformity route, cover reads, dyadic tolerance ledger, Real seal, and structural rows.

The interval-cover computation first reads $I$ and $F$. Thus every point-facing cover read is tied to the displayed located interval and to a listed finite mesh row before any neighborhood witness can be inspected.

The radius and cover checks are finite. The row $L$ supplies the positive radius ledger, while $D$ records the mesh-width and containment comparisons that let $R$ transfer a cover witness from a listed net point to the interval point under inspection.

The sibling uniformity row $C$ is only a compatibility route. It can be consumed by a Heine-Cantor-style downstream reader, but it does not replace the finite-net cover row or introduce an open-cover compactness principle.

Finally, the refused objects have no carrier coordinate. Since $T,Q,P,N$ transport, replay, source, and name only the displayed packet, no accepted $\HeineBorelFiniteNetUp$ certificate exports unrestricted Heine-Borel compactness, a selected subcover, countable choice, host equality, or ambient Real completeness.
\end{proof}

\begin{theorem}[Heine-Borel finite-net interval-cover route]
\label{thm:heine-borel-finite-net-interval-cover-route}
Every interval-cover read from an accepted $\HeineBorelFiniteNetUp$ packet factors through the finite route
$$
\LocatedClosedIntervalFiniteNetUp \longmapsto \LebesgueNumberLemmaUp \longmapsto \DyadicRatCoreUp \longmapsto \RealUp .
$$
The $\HeineCantorUp$ row is available as the sibling compact-uniformity route over the same finite-net and radius data, but the exported certificate remains the displayed finite-net interval-cover route.
\end{theorem}
\begin{proof}
Project the accepted carrier $H=(I,F,L,C,R,D,E,T,Q,P,N)$. The interval source and finite-net row are $I$ and $F$, the positive-radius row is $L$, the dyadic comparison ledger is $D$, and the terminal endpoint seal is $E$.

Replay through $Q$ keeps that order. A cover request opens the finite net, reads the positive-radius ledger, checks the dyadic slack against the mesh and endpoint rows, and only then exposes the Real-facing endpoint seal.

The row $C$ records the compatible Heine-Cantor route for consumers that also need compact uniformity. It remains tied to the same finite rows, so it cannot introduce a selected subcover, an unrestricted open-cover theorem, or an ambient completeness coordinate.
\end{proof}

\begin{theorem}[Heine-Borel finite-net obligation lattice]
\label{thm:heine-borel-finite-net-obligation-lattice}
Any accepted $\HeineBorelFiniteNetUp$ packet must be read in the obligation order
$$
\LocatedClosedIntervalFiniteNetUp \longmapsto \LebesgueNumberLemmaUp \longmapsto \DyadicRatCoreUp \longmapsto \RealUp .
$$
The $\HeineCantorUp$ row is only a sibling uniformity consumer over the same finite interval, positive-radius, dyadic, and Real-seal rows. The packet does not export general open-cover compactness, unrestricted Heine-Borel, selected subcovers, quotient interval representatives, countable choice, host real equality, or ambient $\RealUp$ completeness.
\end{theorem}
\begin{proof}
Project the accepted carrier $H=(I,F,L,C,R,D,E,T,Q,P,N)$. The finite interval source and net rows are $I$ and $F$, the positive-radius row is $L$, the dyadic tolerance ledger is $D$, and the terminal Real-facing seal is $E$.

The cover-read row $R$ can be consumed only after the interval and finite-net rows have been displayed, and the replay row $Q$ keeps the Lebesgue-number row ahead of every dyadic slack check. Thus the finite-net obligation surface is ordered by interval net, positive radius, dyadic comparison, and Real seal.

The row $C$ is read differently. It is available to a $\HeineCantorUp$ consumer that needs compact uniformity over the same finite rows, but it does not precede the finite-net route and does not replace $F,L,D,E$.

Finally the structural rows $T,P,N$ transport, source, and name only the displayed packet. Since the refused objects have no coordinate in $H$, the obligation lattice cannot escape to unrestricted compactness, selected subcovers, quotient representatives, countable choice, host equality, or ambient completeness.
\end{proof}

\begin{theorem}[Heine-Borel finite-net scope witness]
\label{thm:heine-borel-finite-net-scope-witness}
Every accepted $\HeineBorelFiniteNetUp$ packet reads only the displayed located-closed-interval finite-net row, the Lebesgue-number positive-radius row, the Heine-Cantor sibling row, the $\DyadicRatCoreUp$ tolerance ledger, the $\RealUp$ seal, componentwise $\hsame$ transport, displayed $\Cont$ replay, $\Pkg$ provenance, and the local $\NameCert$ row. Thus the chapter scope is closed at the finite-net interval-cover certificate boundary and does not extend to general open-cover compactness, unrestricted Heine-Borel compactness, selected subcovers, quotient interval representatives, countable choice, host real equality, or ambient $\RealUp$ completeness.
\end{theorem}
\begin{proof}
The carrier definition \autoref{def:heine-borel-finite-net-carrier} lists exactly the rows $I,F,L,C,R,D,E,T,Q,P,N$ and states that the classifier compares accepted packets only through those rows.

Apply \autoref{thm:heine-borel-finite-net-namecert-obligations} to identify the local NameCert obligations, \autoref{thm:heine-borel-finite-net-interval-cover-route} to force every interval-cover read through the located interval, finite net, positive radius, dyadic ledger, and Real seal, and \autoref{thm:heine-borel-finite-net-obligation-lattice} to keep the Heine-Cantor row as a sibling uniformity route over the same finite data.

The refused objects have no coordinate in the carrier and are excluded by all three displayed theorem blocks. Therefore the accepted packet reads exactly the stated finite-net scope.
\end{proof}

\closureat{HeineBorelFiniteNetUp}{seedStr}

\begin{closurestatus}{\HeineBorelFiniteNetUp}
  \constructivestory{A $\HeineBorelFiniteNetUp$ packet is generated from a located closed interval source, a LocatedClosedIntervalFiniteNet row, a LebesgueNumberLemma positive-radius ledger, a HeineCantor sibling route, finite cover reads, a DyadicRatCore tolerance ledger, a RealUp seal, componentwise hsame transport, displayed Cont replay, Pkg provenance, and a local NameCert row.}
  \theoryclosure{\scopedClosure}
  \scopeclosed{\autoref{def:heine-borel-finite-net-carrier}, \autoref{thm:heine-borel-finite-net-namecert-obligations}, and \autoref{thm:heine-borel-finite-net-interval-cover-route} seed the constructive finite-net interval-cover certificate over located interval, positive-radius, dyadic tolerance, Heine-Cantor sibling, and Real seal rows.}
  \formalstatus{\unformalizedV}
  \bridgestatus{none}
  \notclaimed{This chapter does not close general open-cover compactness, unrestricted Heine-Borel compactness, selected subcovers, quotient interval representatives, countable choice, host real equality, ambient $\RealUp$ completeness, Lean verification, or bridge checking.}
  \upgradepath{Obligation closure requires checked carrier admission, finite-net interval exposure, radius compatibility, sibling uniformity compatibility, cover-read exactness, non-escape, transport, replay, provenance, and local naming for BEDC.Derived.HeineBorelFiniteNetUp.}
\end{closurestatus}
